\section{Réactivité de l'acide propanoïque}
L'acide propanoïque est un liquide incolore, utilisé en parfumerie pour
synthétiser des composants de parfums, et en médecine vétérinaire pour traiter
les troubles digestifs de certains animaux d'élevage.

Cet exercice vise~:
\begin{itemize}
    \item l'étude d'une solution aqueuse d'acide propanoïque~;
    \item l'étude de la réaction entre l'acide propanoïque et l'éthanol.
\end{itemize}

\subsection{Étude d'une solution aqueuse d'acide propanoïque}
On dispose d'une solution aqueuse ($S_A$) d'acide propanoïque $\ce{C2H5CO2H}$ de
concentration molaire $C_A$ et de volume $V$. La mesure du $\mathrm{pH}$ de la
solution donne la valeur $\mathrm{pH}=3,59$.

L'équation chimique de la réaction entre l'acide propanoïque et l'eau s'écrit~:
\[
    \ce{C2H5CO2H_{(a q)} + H2O$_{(\ell)}$ <=>
    C2H5CO$_2^{-}${(aq)} + H3O^{+}_{(aq)}}
\]

\begin{donnees}
    \begin{itemize}
        \item $\mathrm{pK_A}
              (\ce{C2H5CO2H_{(aq)}\ttslash{}C2H5CO$_2^{-}$_{(aq)}})=4,85$
    \end{itemize}
\end{donnees}
\begin{enumerate}
    \item Donner l'expression de la constante d'acidité $K_{A}$ du couple
          $\ce{C_{2} H_{5} CO2 H_{(aq)}\ttslash{}C2 H5 CO$_2^-$_{(aq)}}$.
          Déduire l'expression du $\mathrm{pH}$ de la solution $(S_{A})$ en
          fonction du $\mathrm{pK_{A}}$ du couple
          $\ce{C2H5CO2H_{(aq)}\ttslash{}C2H5CO$_2^-$_{(aq)}}$ et des
          concentrations $[\ce{C2H5CO2H_{(aq)}}]$ et
          $[\ce{C2H5CO$_2^-$_{(aq)}}]$ en solution.
    \item En exploitant le tableau d'avancement de la réaction, montrer que le
          taux d'avancement final de la réaction s'écrit sous la forme
          $\tau=\frac{1}{1+10^{\mathrm{pK_{A}}-\mathrm{pH}}}$. Calculer sa
          valeur.
    \item Déterminer la valeur de $C_{A}$.
    \item Pour s'assurer de la valeur de $C_{A}$, on dose le volume
          $V_{A}=\qty{20}{\mL}$ de la solution $(S_{A})$ par une solution
          aqueuse $(S_{B})$ d'hydroxyde de sodium
          $\ce{Na^+_{(aq)} + HO^-_{(aq)}}$ de concentration molaire
          $C_{B}=\qty{1e-2}{\mole\per\L}$.
          \begin{enumerate}
              \item Écrire l'équation chimique de la réaction du dosage supposée
                    totale.
              \item Indiquer, en justifiant la réponse, si la solution obtenue à
                    l'équivalence est acide, basique ou neutre.
              \item Le volume de la solution $(S_{B})$ versé pour atteindre
                    l'équivalence acido-basique est $V_{B,E}=\qty{9,8}{\mL}$.
                    Retrouver la valeur de $C_{A}$.
              \item On considère le mélange lorsque le volume de la solution
                    $\left(S_{B}\right)$ versé est $V_{B}=\frac{V_{B, E}}{2}$.
                    \begin{enumerate}
                        \item En exploitant le tableau d’avancement de la
                              réaction du dosage, montrer que
                              $[\ce{C2H5CO2H_{(aq)}}]=
                              \frac{C_B\cdot V_{B,E}}{2\cdot V_{A}+V_{B,E}}$.
                        \item Déduire la valeur du $\mathrm{pH}$ du mélange dans
                              ce cas.
                    \end{enumerate}
          \end{enumerate}
\end{enumerate}

\subsection{Étude de la réaction entre l'acide propanoïque et l'éthanol}
À l'instant $t_{0}=0$, on introduit dans un ballon, $n_{1}=\qty{0,3}{\mole}$
d'acide propanoïque $\ce{C2H5CO2H}$, $n_{2}=\qty{0,3}{\mole}$ d'éthanol
$\ce{C2H5OH}$ et quelques gouttes d'acide sulfurique concentré. On
homogénéise le mélange et on le maintient, durant toute l'expérience, à une
température constante en utilisant le montage de la
figure~\ref{fig:chauffage_a_reflux}. Le volume total du mélange est
$V=\qty{40}{\mL}$.

\begin{figure}[H]
    \centering
    \begin{subfigure}{0.45\textwidth}
        \centering
        \includegraphics[width=\textwidth]{chauffage_a_reflux}
        \caption{}
        \label{fig:chauffage_a_reflux}
    \end{subfigure}
    \hfill
    \begin{subfigure}{0.45\textwidth}
        \centering
        \includegraphics[width=\textwidth]{avancement}
        \caption{}
        \label{fig:avancement}
    \end{subfigure}
    \caption{}
\end{figure}

\begin{enumerate}
    \item Donner le nom du montage de la figure~\ref{fig:chauffage_a_reflux} et
          associer à chaque chiffre le terme correspondant parmi << entrée de
          l'eau - sortie de l'eau - chauffe ballon - réfrigérant à boules -
          mélange réactionnel - ballon - support élévateur >>.
    \item Écrire, en utilisant les formules semi-développées, l'équation
          chimique de la réaction entre l'acide propanoïque et l'éthanol. Nommer
          le composé organique $(E)$ obtenu.
    \item On suit l'évolution de l'avancement $x$ de la réaction en fonction du
          temps, on obtient la courbe de la figure~\ref{fig:avancement}.
          \begin{enumerate}
              \item Déterminer la valeur du temps de demi-réaction $t_{1/2}$.
              \item Déterminer les valeurs de la vitesse volumique de la
                    réaction, en unité (\unit{\mol\per\liter\per\hour}), aux
                    instants $t_{0}=0$ et $t_{1}=\qty{10}{\hour}$. Interpréter
                    qualitativement l'évolution de la vitesse volumique.
          \end{enumerate}
    \item Calculer le rendement $r_{1}$ de la réaction. Comment peut-on
          augmenter ce rendement~?
    \item On peut obtenir le même composé organique (E) à partir de l'éthanol
          et d'un composé organique (A) de formule~:
          \chemfig{C_2H_5-CO-O-CO-C_2H_5}.
          \begin{enumerate}
              \item Préciser le groupe caractéristique du composé $(A)$.
              \item En partant des mêmes quantités de matière
                    $n(A)=n(\text{éthanol})=\qty{0,3}{\mole}$, le rendement de
                    la réaction est $r_{2}$. Comparer, en justifiant, $r_{2}$
                    avec $r_{1}$.
          \end{enumerate}
\end{enumerate}
